\section{Derivations and numerical implementation}
\label{app:math}
\setcounter{table}{0}\setcounter{figure}{0}
\subsection{Posterior allocation and the loss identity}
Under equation~\eqref{eq:model}, $\Cov(\mu,s)=\tau\Sigma$ and $\Var(s)=\tau\Sigma+\Lambda$. The Gaussian conditional mean is therefore $m=\tau\Sigma(\tau\Sigma+\Lambda)^{-1}s$. Multiplying by $\Sigma^{-1}/\gamma$ gives equation~\eqref{eq:heterogeneous}; no commutation of $\Sigma$ and $\Lambda$ is needed. For any position vector $x$, the first-order condition $m=\gamma\Sigma x^H$ implies
\begin{align}
\bar U(x^H\mid s)-\bar U(x\mid s)
&=(x^H-x)'\gamma\Sigma x^H
 -\frac{\gamma}{2}\bigl(x^{H\prime}\Sigma x^H-x'\Sigma x\bigr)\notag\\
&=\frac{\gamma}{2}(x^H-x)'\Sigma(x^H-x).
\end{align}
Taking $x=x^{S,*}$ yields equation~\eqref{eq:loss}. Subtracting and adding $\bar U(x^{S,*}\mid s)$ decomposes the loss of the mean-scalar allocation into compression and choice components. The identity holds draw by draw.

\subsection{When is a scalar adjustment sufficient?}
\label{app:scalar-sufficiency}
The nonnegative loss identity establishes the direction of the oracle comparison. To assess whether the restriction matters economically, we need conditions controlling its magnitude. The following consequences of the same quadratic objective provide such conditions. They hold conditional on the observed signal and the specified risk and uncertainty matrices.

\paragraph{Exact agreement on relevant exposures.}
For any finite $\kappa\geq0$, define the scalar position $q_\kappa=x^S(\kappa)$, the scalar system matrix $B_\kappa=\Sigma+\kappa V$, the penalty difference $D_\kappa=\Lambda/\tau-\kappa V$, and the heterogeneous system matrix $M=\Sigma+\Lambda/\tau=B_\kappa+D_\kappa$. Both system matrices are positive definite. Their first-order conditions give
\begin{equation}
q_\kappa-x^H=B_\kappa^{-1}D_\kappa x^H
             =M^{-1}D_\kappa q_\kappa.
\label{eq:position-gap}
\end{equation}
For example, $B_\kappa q_\kappa=Mx^H=s/\gamma$ implies the first equality; rearranging $Mq_\kappa=s/\gamma+D_\kappa q_\kappa$ gives the second. Consequently, $q_\kappa=x^H$ if and only if $D_\kappa q_\kappa=0$.

In the diagonal uncertainty specification, this condition becomes
\begin{equation}
\sigma_i^2\left(\frac{\lambda}{\tau}h_i-\kappa\right)q_{\kappa,i}=0
\quad\text{for every }i.
\label{eq:exact-scalar}
\end{equation}
Uniform uncertainty gives agreement for every signal at $\kappa=\lambda/\tau$. For a particular signal, agreement also occurs when all nonzero scalar positions have the same uncertainty level matched by $\kappa$. Uncertainty differences on zero positions then leave the allocation unchanged. Small values of $D_\kappa q_\kappa$ give a small position gap when multiplication by $M^{-1}$ does not greatly amplify them. This condition concerns the interaction of uncertainty and exposure, including the risk-matrix geometry.

\paragraph{A sufficient bound on proportional utility loss.}
Let $L_\kappa=\bar U(x^H\mid s)-\bar U(q_\kappa\mid s)$ denote loss at the specified scalar, and let $U_H=\bar U(x^H\mid s)$. At the optimum, $U_H=\gamma x^{H\prime}\Sigma x^H/2$. For $s\ne0$, this quantity is positive. Define the symmetric positive-definite square root $\Sigma^{1/2}$, its inverse $\Sigma^{-1/2}$, the risk-scaled position $y=\Sigma^{1/2}x^H$, and
\begin{equation}
T_\kappa=\Sigma^{1/2}B_\kappa^{-1}D_\kappa\Sigma^{-1/2},
\qquad \eta_\kappa=\|T_\kappa\|_2,
\end{equation}
where the vector norm $\|\cdot\|_2$ is Euclidean and the matrix norm is its induced operator norm. Substituting equation~\eqref{eq:position-gap} into the loss identity yields
\begin{equation}
0\leq\frac{L(s)}{U_H}\leq\frac{L_\kappa}{U_H}
 =\frac{\|T_\kappa y\|_2^2}{\|y\|_2^2}
 \leq\eta_\kappa^2.
\label{eq:scalar-bound}
\end{equation}
The first inequality compares the best scalar with any specified scalar; the last follows from the operator norm. For a chosen proportional-loss tolerance $\epsilon>0$, $\eta_\kappa\leq\sqrt{\epsilon}$ is therefore a sufficient condition for the best scalar to be within $\epsilon$ of heterogeneous utility. The signal-specific middle expression can be much smaller than the uniform upper bound. Large uncertainty differences need not be costly if the held risk-scaled position has little exposure to the directions in which $T_\kappa$ is large. For $s=0$, both allocations and the absolute loss are zero.

\paragraph{Small departures from uniform uncertainty.}
Set $\bar\kappa=\lambda/\tau$, $B=\Sigma+\bar\kappa V$, and $q=B^{-1}s/\gamma$. Fix a diagonal matrix $J=\diag(d_1,\ldots,d_n)$ with $\sum_i d_i=0$. Let $h_i(\theta)=1+\theta d_i$, restricting the scalar perturbation $\theta$ so every $h_i(\theta)>0$. Define $D_0=\bar\kappa V^{1/2}JV^{1/2}$. The heterogeneous matrix is then $B+\theta D_0$. A local inverse expansion at $\theta=0$ gives
\begin{align}
x^H(\theta)-q&=-\theta B^{-1}D_0q+O(\theta^2),\\
L_{\bar\kappa}(\theta)
&=\frac{\gamma\theta^2}{2}
q'D_0B^{-1}\Sigma B^{-1}D_0q+O(\theta^3).
\label{eq:local-loss}
\end{align}
All other inputs, including $s$, remain fixed in this expansion. The position difference is first order, while the utility loss is second order because the objective is locally flat at its optimum. Reoptimizing the scalar can only reduce this loss. Thus modest uncertainty heterogeneity can have a still smaller utility cost. Separately, multiplying the signal by a constant multiplies both positions by that constant and the absolute loss by its square; proportional loss need not fall. These statements distinguish a small dollar-equivalent effect from a small fraction of optimal utility.

\paragraph{Estimated heterogeneity and practical relevance.}
Hold $s$, $\Sigma$ and $\gamma$ fixed, and let $\widehat x^H$ be a feasible position obtained from an estimated positive-semidefinite uncertainty matrix. Evaluate every position using the true conditional objective. Define its estimation penalty as $P=\gamma(\widehat x^H-x^H)'\Sigma(\widehat x^H-x^H)/2$. Completing the quadratic once more gives
\begin{equation}
\bar U(\widehat x^H\mid s)-\bar U(q_\kappa\mid s)=L_\kappa-P.
\label{eq:estimation-tradeoff}
\end{equation}
Estimated heterogeneity improves utility when its estimation penalty is smaller than the scalar restriction cost. The same comparison holds in conditional expectation over an estimator's sampling distribution. This provides a mathematical rationale for retaining a scalar when its approximation loss is small relative to the error introduced by estimating additional uncertainty structure. It does not require the true uncertainties to be identical.

Equations~\eqref{eq:scalar-bound}--\eqref{eq:estimation-tradeoff} concern model utility and a specified economic tolerance. An empirical confidence interval addresses sampling uncertainty in a performance estimate. A confidence interval containing zero does not establish that the bound is small or that the methods are economically equivalent. The appendix supplies conditions under which small differences are plausible; the historical evaluation reports the evidence for incremental net utility under the implemented rule.


\subsection{Scalar search and linear algebra}
The simulation searches in $w=\kappa/(1+\kappa)$, evaluating the zero-position limit at $w=1$ explicitly. For a grid maximum, the neighboring grid points bracket the refinement interval, clipped to $[0,1]$. Golden-section refinement stops when the interval width is at most $10^{-10}$; the refined result and both endpoints are compared. The allocation uses the factor $(1-w)/\gamma$ when represented with the inverse of $(1-w)\Sigma+wV$. The numerical oracle denotes this implemented search over the scalar family. Linear systems are solved directly; a least-squares fallback handles a singular raw benchmark matrix. A constant regression predictor, with centered sum of squares at most $10^{-14}$, produces an intercept-only fit. Identical bootstrap predictions give exactly zero estimated variance.

\subsection{Principal-component statistics}
Suppress the formation-date subscript. For a risk matrix $M$, write $R_M=V^{-1/2}MV^{-1/2}$. Let $d_k$ and $p_k$ denote its ordered eigenvalues and corresponding unit eigenvectors, for $k=1,\ldots,n$. The first three statistics below are computed with $M=\widehat\Sigma$ for raw mean--variance, $M=Q$ for corrected mean--variance and $M=(1-w)Q+wV$ at the selected $w$ for EPO. Because $R_Q=0.95R_{\widehat\Sigma}+0.05I$ and $R_{(1-w)Q+wV}=(1-w)R_Q+wI$, the three matrices share eigenvectors, and their eigenvalues are increasing transformations of one another, so the ordering and groups coincide. The realized component return is $a_k=p_k'V^{-1/2}r^e_{t+1}$, and the projected signal is $b_k=p_k'V^{-1/2}z_t$. Component $k$ belongs to group $1+\lfloor5(k-1)/n\rfloor$. The group risk ratio is the square root of the time average of within-group means of $a_k^2/d_k$. It uses an uncentered second moment. The signal-implied statistic averages $|b_k|/\sqrt{d_k}$ within a group and divides by the same month's all-component mean. Realized performance uses the annualized Sharpe ratio of the monthly group mean of $\operatorname{sign}(b_k)a_k/\sqrt{d_k}$. For positions $x$ of any of the three portfolios, component $k$ contributes $d_k(p_k'V^{1/2}x)^2$ to predicted variance, with $d_k$ taken from $R_Q$; group shares are normalized each month before averaging.

\section{Simulation and estimator-calibration details}
\label{app:simulation}
\setcounter{table}{0}\setcounter{figure}{0}
\subsection{Profile generation and Monte Carlo summaries}
Profiles follow equation~\eqref{eq:simulation-profile} and are drawn before the common return and signal-error innovations. Each profile remains fixed across signal draws; equation~\eqref{eq:simulation-summary} defines the reported mean and proportional losses. The simulation seed is 20261023. A mean-loss standard error is the sample standard deviation across signal draws divided by the square root of the draw count; the interval is the mean plus or minus 1.96 standard errors. Means, standard errors and interval endpoints are multiplied by 12 for annual reporting, leaving loss shares unchanged. Alignment groups use a stable ascending sort and equal-count splitting. Table~\ref{tab:simulation} reports both actual and target dispersions.
\input{generated/table_simulation}

\subsection{Controlled uncertainty-estimation error}
The estimation-error experiment reconstructs the same fixed true uncertainty vectors and 1,000 signal draws using seed 20261023. Independent standard-normal estimation errors use seed 20261112. Equation~\eqref{eq:controlled-h-error} defines their multiplicative effect at $a=0$, $0.5$ and $1$. Normalization fixes the mean estimated relative uncertainty at one in every draw; the experiment therefore varies relative estimation error at a known common scale. At $a=0$, the estimates equal the true values. The error parameter describes the synthetic estimator used in this experiment.

For each draw, the researcher evaluates all positions with the same true conditional mean and covariance. Mean losses are arithmetically annualized by multiplying by $12\times10^4$. Monte Carlo standard errors are sample standard deviations across independent paired draws divided by $\sqrt{1{,}000}$; interval endpoints are the mean plus or minus 1.96 standard errors. Ratios divide mean losses by mean oracle utility. Table~\ref{tab:estimation-tradeoff} reports the utility differences and the frequency of improvement. The specified mean-scalar cost $L_{\bar\kappa}$ includes the choice loss in Table~\ref{tab:simulation}, whereas that table's compression column uses the best scalar.
\input{generated/table_estimation_tradeoff}

\subsection{Calibration data-generating process}
Let $e_{i,k+1}$ be standard-normal return innovations, independent over time and correlated 0.30 across industries. Let $\xi_{i,k+1}$ be independent standard normals, independent of all return innovations. The synthetic predictor starts at $z_{i,0}=s_z v_{i,0}$, where $s_z=0.056$ and the $v_{i,0}$ are independent standard normals. With persistence $\phi=0.92$ and feedback $f\in\{0,0.4\}$, generate
\begin{align}
z_{i,k+1}&=\phi z_{i,k}+s_z\sqrt{1-\phi^2}
       \left(fe_{i,k+1}+\sqrt{1-f^2}\xi_{i,k+1}\right),\\
r^e_{i,k+1}&=a+bz_{i,k}+\nu_i e_{i,k+1},
\end{align}
where $a=0.005$, $b=0.10$, and $\nu_i$ is equally spaced between 0.04 and 0.08 across the 49 industries. Each history provides pairs $k=0,\ldots,120$. We train on the first 120 and forecast using $z_{i,120}$, whose true next-return conditional mean is $a+bz_{i,120}$. Current predictors remain independent of the next return innovation in both feedback designs. Relative uncertainty uses the sample variance of the latest 60 training responses, matching the empirical estimator.

For zero feedback, define $\bar z_i$ as the mean training predictor and $S_{zz,i}=\sum_{k=0}^{119}(z_{i,k}-\bar z_i)^2$. The exact conditional variance of the fitted-mean error is
\begin{equation}
v_i^{\mathrm{reg}}=\nu_i^2\left[\frac{1}{120}
 +\frac{(z_{i,120}-\bar z_i)^2}{S_{zz,i}}\right].
\end{equation}
With feedback, conditioning on the complete predictor history changes the relationship between predictors and training residuals. We retain this formula as an analytical reference, and assess actual calibration using errors against the known mean. The rank reference for relative uncertainty is $v_i^{\mathrm{reg}}/(b^2s_z^2+\nu_i^2)$.

The calibration seed is 20261101, with separate child streams for histories and bootstrap sampling. Innovations and bootstrap dates are paired across feedback designs. Coverage is computed for normal intervals centered on the fitted mean with standard error $\sqrt{u_i}$, and for empirical percentile intervals of bootstrap predictions, at nominal 90\% and 95\% levels. We evaluate 49,000 industry-history observations per design. Mean-statistic standard errors aggregate centered observations within each independent history before squaring and summing; ratio standard errors apply the same calculation to the linearized ratio. Rank-correlation standard errors use the dispersion across histories. Table~\ref{tab:calibration} reports the aggregate quantities.
\input{generated/table_calibration}
\FloatBarrier

\section{Data construction and replication discrepancies}
\label{app:data}
\setcounter{table}{0}\setcounter{figure}{0}
The source archives were downloaded on October 23, 2026 and identify the CRSP vintage as August 2026. The industry returns follow the French Data Library's industry assignments and value-weighting convention \citep{frenchdata}. Original archives, processed returns and the source manifest accompany the research code. The benchmark contains 924 calendar months from January 1942 through December 2018. The file has 1,656 missing industry-months in that interval, many of which are excluded by formation-date eligibility. One missing realized return affects an industry held by the benchmark: Rubbr in July 1943. That portfolio month remains missing. The January selection windows for 1944--1948 therefore fail the complete-history requirement, removing another 60 months. The common evaluation sample has 863 months. Eligible industries range from 42 to 49, with mean 46.8 across formation dates. All 49 remain eligible from June 1974 onward.

\subsection{Accounting for differences from the published values}
The primary replication follows the conventions in Section~\ref{sec:method}. The description of monthly expanding-history selection comes from \citet[p.~138]{pedersen2021}, in ``Global Asset Classes: Beating Time-Series Momentum,'' under ``Performance of EPO vs. benchmark portfolios.'' The text describes monthly parameter choice using an expanding history; the Figure~2 note on the same page specifies maximizing past Sharpe ratio over a finite grid for use in the next period. For equities, the note to Table~5 (p.~142) states that selection uses only past data, without restating the updating frequency or window. Our diagnostic applies the monthly expanding rule to Equity~1 as an interpretation of that procedure.

Table~\ref{tab:conventions} isolates a sequence of convention changes on the current data vintage. The diagnostic uses the published grid, omits monthly volatility targeting and sets the single held missing return to zero, producing a Sharpe ratio of 0.7971 over all 924 months. Restricting evaluation to the common evaluation months raises it to 0.8238. Changing to the seven-value grid in Section~\ref{sec:method} leaves that rounded result unchanged. Scaling positions to 10\% predicted volatility produces 0.8311, and changing to January selection from the preceding 60 months produces 0.7703.

This accounting attributes the largest downward step along the chosen path to the selection-window convention. It does not identify a unique additive cause for every difference from the published result: the sequence is order-dependent, and the historical data vintage also differs. The zero-filled calculation is confined to this diagnostic exercise. The primary result retains missing held returns and uses the common sample described above. The full convention cross-comparison is available in the machine-readable benchmark results.

The two fixed portfolios with the largest differences in Table~\ref{tab:replication}, raw mean--variance and EPO with $w=0.99$, follow the same path without the selection steps (Table~\ref{tab:conventions-fixed}). Applying the published procedure to the current data, with the missing held return set to zero, gives 0.2195 and 0.7535, within 0.03 of the published 0.19 and 0.73. Restricting evaluation to the common months raises them to 0.2336 and 0.7904. Monthly scaling to 10\% predicted volatility under $Q_t$ raises them to 0.3221 and 0.8313. The raw portfolio takes its direction from $\widehat\Sigma_t$ but its scale from $Q_t$. Because $Q_t$ assigns more risk than $\widehat\Sigma_t$ to the small-eigenvalue bets that dominate this portfolio, months with the most extreme bets receive the smallest positions. Along this path, the scaling step accounts for most of the raw portfolio's difference. For $w=0.99$, the common sample and the scaling contribute about 0.04 each.
\input{generated/table_conventions}
\input{generated/table_conventions_fixed}
\reportfigure[H]{sample_overview.pdf}{Benchmark coverage and industry return distributions. The left panel shows formation-date eligibility and marks the missing held return and the skipped selection years. The other panels summarize each industry's eligible excess returns within the common evaluation months. Mean returns and volatility are annualized.}{fig:benchmark-sample}{1}
\reportfigure[H]{selected_w.pdf}{Annual benchmark penalty selection. Each January's choice maximizes gross Sharpe ratio over the preceding 60 complete months. Monthly positions are scaled to 10\% predicted volatility. The vertical spacing represents grid categories rather than equal numerical distances; shaded years have no admissible selection window.}{fig:benchmark-w}{1}
\reportfigure[H]{sample_coverage.pdf}{Coverage of the empirical extension. The left panel compares observed industry returns with forecast-eligible industries. The right panel gives annualized volatility distributions using each industry's eligible observations in each evaluation period. The earlier period contains 28,104 eligible industry-months and the later period 4,116.}{fig:sample}{1}
\FloatBarrier

\section{Forecast and relative-uncertainty diagnostics}
\label{app:forecast}
\setcounter{table}{0}\setcounter{figure}{0}
Monthly forecast deciles use a stable ascending sort and near-equal counts. The plots report equal-weight mean forecasts and next-month realized excess returns in each group. The predictive $R^2$ in Table~\ref{tab:forecast} is one minus total squared forecast error divided by total squared error of the rolling historical-mean forecast, pooled over eligible industry-months. The historical mean uses the same 120 training responses. The corresponding 95\% intervals are $[-0.0135,-0.0041]$ and $[-0.0200,-0.0047]$ in the two periods. Rank correlations are calculated within a month and then averaged over months.

The uncertainty profile is normalized monthly, so cross-sectional means equal one by construction. Dispersion uses the population standard deviation within a month. Persistence compares ranks on the common industry intersection at the two dates, requiring at least three industries and nonconstant ranks. Monthly diagnostics also compare uncertainty ranks with absolute forecast and asset-volatility ranks. The empirical uncertainty-estimator seed is 20261029, combined with formation month in year-month integer form; evaluation resampling uses seed 20261105. Diagnostic confidence intervals use whole-month blocks, with seed 20261105 combined with a diagnostic identifier and the starting year of the evaluation period.
\reportfigure[H]{forecast_quality.pdf}{Historical forecast diagnostics. Upper panels compare forecast and realized mean excess returns by forecast decile, before costs, with 95\% whole-month moving-block intervals for realized means. The lower panel is the trailing 60-month average cross-sectional Spearman correlation between forecasts and realized returns. Intervals use 1,000 draws with 12-month blocks and are pointwise.}{fig:forecast}{1}
\reportfigure[H]{reliability_diagnostics.pdf}{Dispersion and persistence of estimated relative uncertainty. The upper panel shows the monthly coefficient of variation. The middle panel compares uncertainty ranks at one- and twelve-month lags. The lower panel compares uncertainty ranks with absolute forecast and asset-volatility ranks. The bottom two panels display trailing twelve-month averages. Larger uncertainty values imply a stronger relative penalty.}{fig:reliability}{1}
\reportfigure[H]{reliability_heatmap.pdf}{Annual mean relative uncertainty by industry. Each cell is the arithmetic mean of monthly normalized uncertainty values. The color scale is logarithmic and saturates outside 0.25 to 4; gray denotes unavailable observations. Industry coverage follows the primary forecast eligibility rule.}{fig:heatmap}{.94}
\FloatBarrier

\section{Selection and portfolio transmission}
\label{app:holdings}
\setcounter{table}{0}\setcounter{figure}{0}
All candidate accounts start in January 1960, using available historical forecasts. The unpenalized candidate becomes insolvent in June 1960 at target gross exposure 222.3, and the $w=0.25$ candidate in December 1961 at target gross exposure 71.0, measured per unit of wealth. Both remain excluded from later selection. These extreme positions arise in the unconstrained candidate allocation. At each annual decision, a candidate must still be alive and have 60 complete prior monthly net returns. Figure~\ref{fig:selection} shows selected values and the eligible historical scores across the 56 evaluation years. The candidate-exit and selection records are retained separately from the displayed aggregate results.

Figure~\ref{fig:transmission} shows gross exposure and predicted volatility for the primary portfolios. The composition measures defined in Section~\ref{sec:method} appear in Figure~\ref{fig:composition}. Holdings include entering and exiting industries, with zero target positions for liquidations.
\reportfigure[H]{selection_diagnostics.pdf}{Annual common-penalty selection in the empirical extension. The upper panel shows the selected penalty. The lower panel shows annualized net certainty equivalents over the preceding 60 months for eligible scalar candidates. Color limits use the fifth and ninety-fifth percentiles of displayed scores; excluded candidates are omitted. Both strategies use the selected scalar penalty.}{fig:selection}{1}
\reportfigure[H]{portfolio_transmission.pdf}{Gross exposure and annual predicted volatility under the common risk matrix. These are the primary fixed-risk-aversion portfolios before volatility targeting and capping. Figure~\ref{fig:composition} reports their composition distance and risk-direction correlation.}{fig:transmission}{1}
\FloatBarrier

\section{Performance paths and sensitivity intervals}
\label{app:performance}
\setcounter{table}{0}\setcounter{figure}{0}
Figure~\ref{fig:equity} displays total net wealth and drawdown for the primary strategies. Each period is rebased to initial wealth one for presentation; portfolio positions continue through the boundary. Table~\ref{tab:intervals} and Figure~\ref{fig:intervals} give paired uncertainty intervals under alternative return-block lengths. Figure~\ref{fig:robust} reports the cost and entry-rule intervals using twelve-month blocks. Table~\ref{tab:performance} reports the 12-month-block intervals for the common position-sizing comparison. The bootstrap uses seed 20261105 combined with the period start year and block length, identically across position-sizing rules.
\input{generated/table_intervals}
\reportfigure{equity_curves.pdf}{Total net wealth and drawdown of the primary matched portfolios. Returns include risk-free cash financing, 10 basis points per dollar traded and a 50-basis-point annual short fee. Each display period is rebased independently; underlying holdings and cost accounting continue in January 2019.}{fig:equity}{1}
\reportfigure{delta_ce_intervals.pdf}{Block-length sensitivity of the heterogeneous-minus-scalar annual net certainty equivalent difference. Each horizontal interval is the 2.5th to 97.5th percentile of 1,000 paired moving-block resamples of realized net returns. The central marker is the original-sample estimate, and the dashed line marks zero. Units are annual basis points.}{fig:intervals}{1}
\reportfigure{robustness.pdf}{Cost and admission sensitivity with paired uncertainty intervals. Cost scenarios retain the primary penalty sequence and recompute position drift; the annual short fee remains 50 basis points. Requiring 60 forecast-ready months reruns candidate histories and selection. Intervals use 1,000 paired twelve-month block resamples and are expressed in annual utility basis points.}{fig:robust}{1}
\FloatBarrier

\subsection{A fixed-label diagnostic for estimated uncertainty}
\label{app:label-placebo}
The label-assignment diagnostic assesses the usefulness of matching estimated uncertainty to individual industries. Starting from the forecasts and normalized uncertainty profiles, we draw 100 random rankings of the 49 industry labels with seed 20261111. For each ranking, at each formation date we restrict that ranking to the eligible industries and assign the corresponding permutation of $\widehat h_t$ to the canonical eligible-industry order. The numerical distribution of $\widehat h_t$ is preserved exactly each month. The mapping remains fixed whenever the eligible universe is unchanged; it adjusts mechanically at an eligibility change. This avoids generating independent new label assignments every month.

Each shuffled portfolio uses the original forecasts, covariance matrix, selected scalar penalty, fee rates and evaluation dates. We evaluate both fixed risk aversion and the predicted-volatility target with gross cap. Every account starts from cash in January 1970 and carries its own positions and costs continuously through 2025. All 100 paths remain solvent under both policies. No tuning parameter is chosen using the diagnostic outcomes.
Table~\ref{tab:label-placebo} presents the resulting utility comparisons.

Before trading and short costs, the original matching exceeds 99 and 77 of the 100 shuffled portfolios in the earlier period under fixed risk aversion and targeted exposure, respectively; these counts fall to 19 and 21 in the later period. Costs alone therefore do not account for the change in relative ordering reported in Section~\ref{sec:label-diagnostic}.

The rank summaries describe the 100 specified alternative assignments. For inference on their mean utility benchmark, equation~\eqref{eq:assignment-advantage} defines $D$ as the original strategy's CE minus the average of the 100 individual CEs. The same 1,000 bootstrap date sequences are applied to every strategy within each period, using overlapping blocks of 12 consecutive months. Dates are sampled jointly across both position-sizing rules as well. We recompute each strategy's CE and their mean benchmark in every sample. The 2.5th and 97.5th percentiles give the interval in Table~\ref{tab:label-placebo}.

These intervals describe temporal sampling uncertainty conditional on the fitted forecasts, selected penalties and the specified assignments. They use the moving-block bootstrap approximation to dependence in monthly strategy returns. The assignment ranks themselves remain descriptive; industry labels are not assumed exchangeable and no permutation significance level is assigned to them. Table~\ref{tab:assignment-distribution} preserves the mean and median summaries of the random strategies.
\input{generated/table_assignment_distribution}
\FloatBarrier
